\documentclass[11pt,a4paper]{article}

\usepackage[margin=2.5cm]{geometry}
\usepackage{graphicx}
\usepackage{booktabs}
\usepackage{amsmath}
\usepackage{hyperref}
\usepackage{float}
\usepackage{subcaption}

\title{Flexible Cable Simulation with Rigid Connectors in IsaacSim}
\author{Ali Bani Asad}
\date{\today}

\begin{document}

\maketitle

% ------------------------------------------------------------------
\section{Introduction}
% ------------------------------------------------------------------

This report presents a method for modeling and simulating a flexible cable
connected to rigid connectors within NVIDIA IsaacSim. Since IsaacSim does
not natively support deformable 1D cables, the cable is approximated using
a chain of rigid capsule segments connected by constrained joints,
providing realistic bending and torsional behavior.


% ------------------------------------------------------------------
\section{Approach}
% ------------------------------------------------------------------

The cable is modeled as a chain of \textbf{60 capsule-shaped rigid bodies}
connected by \textbf{D6 joints} (6 degree-of-freedom joints). Each D6 joint
locks all three translational axes and constrains the three rotational axes
with cone and twist limits, producing cable-like flexibility.

Two rigid connectors are attached to the cable:
\begin{itemize}
    \item \textbf{Top connector} (blue cube): A rigid block fixed to the
          world frame via a fixed joint. The cable's first link attaches to
          its bottom face through a D6 joint with limited rotation.
    \item \textbf{Bottom connector} (red cube): A weighted rigid block
          attached to the cable's free end via a fixed joint, representing
          a plug or end-fitting.
\end{itemize}

A \textbf{green obstacle} (fixed cube) is placed in the scene so the cable
can drape over it and demonstrate realistic contact interaction.


% ------------------------------------------------------------------
\section{Simulation Parameters}
% ------------------------------------------------------------------

\begin{table}[H]
\centering
\caption{Cable and simulation parameters.}
\begin{tabular}{@{}llr@{}}
\toprule
\textbf{Parameter}         & \textbf{Symbol}    & \textbf{Value}  \\
\midrule
Number of links            &  $N$               & 60              \\
Total cable length         &  $L$               & 1.0 m           \\
Link radius                &  $r$               & 4.0 mm          \\
Link height (cylinder)     &  $h$               & 8.67 mm         \\
Segment spacing            &  $\Delta s$        & 16.67 mm        \\
Total cable mass           &  $m_{\text{total}}$ & 1.0 kg         \\
Per-link mass              &  $m_i$             & 16.67 g         \\
Cone limit (swing)         &  $\theta_c$        & $8.0^\circ$     \\
Twist limit                &  $\theta_t$        & $5.8^\circ$     \\
Linear damping             &  $d_l$             & 0.2             \\
Angular damping            &  $d_a$             & 1.0             \\
Joint damping              &  $d_j$             & 0.05            \\
Physics time step          &  $\Delta t$        & 1/240 s         \\
Solver position iterations &                    & 64              \\
Solver velocity iterations &                    & 8               \\
Solver type                &                    & TGS             \\
CCD                        &                    & enabled         \\
\bottomrule
\end{tabular}
\end{table}

\begin{table}[H]
\centering
\caption{Rigid connector parameters.}
\begin{tabular}{@{}lrr@{}}
\toprule
                    & \textbf{Top connector} & \textbf{Bottom connector} \\
\midrule
Size (cube side)    & 6 cm                   & 4 cm                      \\
Mass                & 5.0 kg                 & 0.5 kg                    \\
Constraint          & Fixed to world         & Free (weighted)           \\
\bottomrule
\end{tabular}
\end{table}


% ------------------------------------------------------------------
\section{Results}
% ------------------------------------------------------------------

The simulation was run for 10 seconds. Key frames captured during the
simulation are shown in Figure~\ref{fig:keyframes}. The cable is given an
initial lateral velocity of 1.5 m/s at the bottom connector to induce
swinging motion toward the obstacle.

\begin{figure}[H]
\centering
\begin{subfigure}[b]{0.48\textwidth}
    \includegraphics[width=\textwidth]{cable_output/frame_t0s.png}
    \caption{$t = 0$ s --- initial configuration}
\end{subfigure}
\hfill
\begin{subfigure}[b]{0.48\textwidth}
    \includegraphics[width=\textwidth]{cable_output/frame_t2s.png}
    \caption{$t = 2$ s --- cable swinging}
\end{subfigure}

\vspace{0.5cm}

\begin{subfigure}[b]{0.48\textwidth}
    \includegraphics[width=\textwidth]{cable_output/frame_t5s.png}
    \caption{$t = 5$ s --- contact with obstacle}
\end{subfigure}
\hfill
\begin{subfigure}[b]{0.48\textwidth}
    \includegraphics[width=\textwidth]{cable_output/frame_t10s.png}
    \caption{$t = 10$ s --- settling}
\end{subfigure}
\caption{Key frames from the cable simulation.}
\label{fig:keyframes}
\end{figure}

A full 10-second video recording of the simulation is available in the
supplementary file \texttt{cable\_simulation.mp4}.


% ------------------------------------------------------------------
\section{Challenges}
% ------------------------------------------------------------------

\begin{itemize}
    \item \textbf{Stability}: Long rigid-body chains are prone to jitter.
          High solver iteration counts (64 position / 8 velocity) and the
          TGS solver were required for stable behavior.
    \item \textbf{Thin-body collisions}: The small capsule radius (4 mm)
          required continuous collision detection (CCD) to prevent
          tunneling at high velocities.
    \item \textbf{Damping tuning}: Balancing linear, angular, and joint
          damping was critical --- too little causes a frictionless pendulum
          chain; too much makes the cable overly rigid.
    \item \textbf{Joint limits}: The cone and twist limits control bending
          stiffness. Smaller angles produce a stiffer cable but can cause
          constraint violations if set too tight.
\end{itemize}


% ------------------------------------------------------------------
\section{Conclusion and Next Steps}
% ------------------------------------------------------------------

The capsule-chain approach successfully models a flexible cable connected
to rigid connectors in IsaacSim, with realistic gravity, contact, and
bending behavior. Possible next steps include:

\begin{itemize}
    \item Integrating the cable with a robot manipulator for cable
          manipulation tasks.
    \item Exploring IsaacSim's deformable body API as an alternative to
          the rigid-body chain approach.
    \item Adding friction and material properties for more realistic
          surface interactions.
    \item Parameterizing the cable for different lengths, stiffnesses,
          and diameters.
\end{itemize}

\end{document}
